Tool-using agents often invoke a model at every step, even along a repeated read-only
path. Replacing repeated reasoning with deterministic execution could reduce inference
cost, but repetition alone does not make this safe: tool arguments may depend on decisions
not recoverable from observable state, apparently read-only operations may have hidden
effects, and reproducing the same tool calls may still change the agent's final answer.
Recurrence identifies an opportunity for optimization, not permission to remove reasoning.

We introduce \emph{guarded agentic compaction} (\method), a compile-or-retire approach
that converts recurrent execution paths into deterministic programs only when the evidence
supports it. \method grounds every tool argument in the task input or a prior tool result,
treats writes, approvals, handoffs, and unknown effects as hard barriers, protects
compiled execution with input, output, and environment guards, and requires a
finite-sample risk test before deployment. Admission depends jointly on a configured risk
bound and the statistical confidence supported by the available traces. If any requirement
fails, \method keeps the largest justified prefix when possible, or falls back to the
unchanged agent.

Across three live GitHub workflow families and 90 unseen test cases, \method satisfies the
exact task requirements on all 90 cases, compared with 89 of 90 for the unchanged agent,
while reducing model requests by 66.6\%, tokens by 63.1\%, latency by 64.2\%, and estimated
cost by 58.7\%. A time-forward evaluation across five more repositories shows \method
specializing four and refusing the fifth. Experiments on four public trace benchmarks
further show that recurrence and successful replay are not sufficient for specialization.
Under \method's default 0.05 risk bound and associated confidence requirement, NESTFUL,
API-Bank, and BFCL do not provide sufficient qualifying recurrent evidence to establish
that substitution risk is within the admissible bound, so \method conservatively declines
specialization; relaxing the risk or confidence requirements can admit additional
opportunities, making the trade-off between optimization coverage and substitution
assurance explicit. The fourth benchmark demonstrates a distinct limitation: even when
trace evidence supports compilation, the resulting deterministic program can remain
incompatible with the target agent architecture.

\method reframes optimization as evidence-gated specialization: compile only what is
justified, and preserve the agent everywhere else.
